\scanpage{103}
% archive foliation: 102
\textcircled{3} \uline{Corollaire} $X$ schéma normal \deleted{\unsure{et}} loc. \unsure{intègre}

$\unsure{D_i}$ diviseurs\inserted{de Cartier} intègres \deleted{\unsure{réunis}}

\margin{\diagram{\begin{tikzcd}[column sep=small,row sep=small]
K \arrow[r,dash] & \widetilde K \arrow[r,dash] & M \\
X & \widetilde X & X'
\end{tikzcd}}}

$X'$ revêtt \deleted{\unsure{plat}} normal galoisien, groupe $G$, tel que les grpes d'inertie\inserted{\deleted{$\unsure{I_i}$}$\bar I_i$ (d'ordre $n_i$)} \deleted{relatifs à \unsure{chacun des}} $D_i$ \unsure{soient} invariants \deleted{et commutent deux à deux, et que}\inserted{et que} $X'$ soit modérément ramifié p.r. \deleted{aux \unsure{points de}}\inserted{à} $D$. Donc on a, \deleted{où $\xi$ est un pt géom. de $X$}

\deleted{Alors}
\diagram{\begin{tikzcd}[column sep=small]
1 \arrow[r,dash] & \bar I \arrow[r,dash] & G \arrow[r,dash] & \Gamma \arrow[r,dash] & \unsure{1}\mathbin{-}1
\end{tikzcd}}
\deleted{$\bar I$ est \unsure{isomorphe au produit des} $\unsure{\mu(n_i)}$}. De plus,

$X'/\bar I\unsure{{}_a}=\widetilde X$ est \deleted{\unsure{donc}} un rec. \uline{principal} de $X$ de groupe $\Gamma$. Donc la classe $\alpha\in H^2(\Gamma,\bar I)$ de l'extension définit une classe \unsure{de} $\alpha\in H^2(X,\widetilde X\mathbin{\times}^{\Gamma}\bar I)$, \deleted{\unsure{D'}}\inserted{et à fortiori} en $H^2(D,\widetilde D\mathbin{\times}^{\Gamma}\bar I)$. \unsure{Or} on a un isomorphisme canonique
\[
\deleted{\unsure{I_{\widetilde X}}}\widetilde D\mathbin{\times}^{\Gamma}I
\simeq\deleted{\prod\mu_{n_i,D}}\mu_{\unsure{n},D}
\]
i.e. un isom. compatible avec les opérations de $\Gamma$ \deleted{\unsure{et}} sur les faisceaux envisagés sur $(\widetilde X,\Gamma)$
\[
\unsure{\text{en}}\qquad \bar I_{\widetilde D}\deleted{\unsure{\simeq\mu_{\widetilde X}}}\simeq\mu_{\underline n,\widetilde D}
\]
\deleted{\unsure{Pour définir}} cet isomorphisme est tel que, \unsure{en le} \unsure{par fibres}, opérateurs $\bar I_{\widetilde X}\simeq\mu_{\underline n}$, $X'$ est \deleted{\unsure{pareil}} Kummérien sur $\widetilde X$, relatif aux $\widetilde D_i$, à exposants $\underline n=(n_i)$.

Ceci dit, l'image de $\alpha$ dans $H^2(D,\mu_{\underline n})$ est égale à la classe déduite sur $D$ par $c(\underline{O}_X(D))\in H^2(X,\mu_{\underline n})$, i.e.
\margin{\fbox{\parbox{0.8\linewidth}{\deleted{On peut supposer} $X$ intègre. Le choix d'un pt géométrique $\bar\xi_i$ de $X'$ \unsure{au-dessus d'un pt gén.} $D_i$ définit un isom $I_i\simeq\mu$.}}}
\diagram{\begin{tikzcd}[column sep=large]
c(T_{D/X})\in H^2(D,\mu_{\underline n}) \arrow[r,bend right=15] & \text{\uline{fibré normal}}
\end{tikzcd}}
